\chapter{Fireballs, Frames, and Four-Vectors}

We begin with a question left over from the discussion of high-energy
collisions. A collision of two heavy nuclei creates a violently changing
object, not a quiet astronomical black hole. What, then, can it mean to
describe the collision by a horizon? The answer leads through a
five-dimensional picture of strongly coupled matter, a set of practical
questions about accelerators, and finally back to the geometry of special
relativity. By the end, velocity and momentum have become four-vectors, and
energy has taken its place beside the three components of momentum.

\section{A Fireball in a Fifth Dimension}

\begin{classroomqa}
\audiencequestion{The object produced in the collision is unstable and
changes shape. How can such a complicated object be called a black hole?}
\lecturerresponse{Picture a fifth spatial direction running between a floor
and a ceiling. Gravity in that direction pulls excitations toward the floor.
The collision deposits a hot lump of energy there; the lump spreads and
changes shape, but it can still possess a horizon, temperature, and entropy.}
\end{classroomqa}

Imagine particles, or more pictorially little strings, moving in this
five-dimensional world. Excitations that fall to the floor appear in four
dimensions as hadrons: protons, neutrons, mesons, and other extended composite
objects. Excitations attached near the ceiling appear more elementary. The
vertical location is therefore a geometric way of encoding scale: the floor
represents the long-distance hadronic regime, while the ceiling represents
short-distance degrees of freedom.

When two hadrons collide, two objects near the floor strike one another and
make a large splash of energy. The result is not a rigid body. It is a hot
horizon-bearing lump, a kind of puddle in the five-dimensional gravitational
field. It is not an astronomical black hole in ordinary four-dimensional
space, nor a Planck-scale black hole attached to the ceiling. It is the
gravitational description of the strongly interacting fireball near the
floor.

At high energy the incoming nuclei are Lorentz-contracted into thin pancakes.
They collide, deposit their energy in a small region, and produce a fluid that
spreads toward equilibrium. Its small viscosity was one of the surprising
features of the quark--gluon plasma. In the five-dimensional description the
same evolution is the relaxation of a distorted horizon.

\begin{figure}[H]
\centering
\includegraphics[width=0.82\textwidth]{lecture_05_figure_01.png}
\caption{The floor-and-ceiling picture of the fifth dimension. The central
horizon-bearing lump spreads along the floor after the collision.}
\lectureframe{00:04:22}
\end{figure}

A horizon does not remain forever. Hawking radiation carries energy away, and
in the four-dimensional description that radiation is the shower of ordinary
hadrons emerging from the collision. The thousands of tracks in a detector
event are, in this language, the particles boiled off the dual black hole.

\subsection{Is the extra dimension real?}

\begin{classroomqa}
\audiencequestion{Is the fifth dimension a spatial dimension, and is the line
drawn between floor and ceiling literally extending in that direction?}
\lecturerresponse{Yes. It is spacelike, and the surface drawn in that
five-dimensional geometry is a horizon. Whether one calls the extra direction
``real'' or regards it as an equivalent mathematical description is a deeper
question.}
\end{classroomqa}

The useful statement is not that gravity replaces quantum chromodynamics.
Rather, two languages may describe the same physics. In the four-dimensional
language there is a hot fireball of quarks and gluons, with no ordinary
gravitational field. In the five-dimensional language there is gravity and a
horizon. If the correspondence is correct, an honest QCD calculation and the
dual gravitational calculation must agree wherever both can be performed.

\begin{classroomqa}
\audiencequestion{QCD has neither gravity nor horizons. Is the black-hole
picture adding new physics, or does it merely redescribe the fireball?}
\lecturerresponse{In four dimensions it is a plasma of quarks and gluons; in
the five-dimensional incarnation it has a gravitational description with a
horizon. The advantage is calculational. Strongly coupled QCD is extremely
difficult, while some horizon properties---the viscosity in particular---are
comparatively easy to calculate.}
\end{classroomqa}

The distinction between ``real dimension'' and ``mathematical analogy'' need
not be settled before the description can be useful. A structure earns a
physical interpretation when it joins the other variables into a simpler,
more coherent set of laws. Still, one should not overstate the case:
applications of gauge/gravity duality to the observed quark--gluon plasma are
models of strongly coupled dynamics, not a proof that real-world QCD has a
known exact gravitational dual.\editorialnote{This qualification makes
explicit the modern status of the QCD application; it does not alter the
five-dimensional mechanism developed here.}

\section{Why Colliders Collide Beams}

The questions now become practical. Why use two moving beams instead of a
single beam and a stationary target? Why are circular accelerators so large?
Why are electrons and protons handled differently? All of these questions are
questions about relativistic energy.

\begin{classroomqa}
\audiencequestion{Why collide two beams rather than fire one energetic beam
into a stationary target?}
\lecturerresponse{What creates new particles is the energy available in the
center-of-mass frame. Two equal and opposite relativistic beams have no net
momentum, so essentially all their energy is available to the collision. A
fixed target forces much of the incident energy to remain as motion of the
center of mass.}
\end{classroomqa}

For two ultrarelativistic particles with opposite momenta of magnitude \(p\),
the lecture-level estimate in \(c=1\) units is
\begin{equation}
E_{\mathrm{cm}}^{\mathrm{collider}}\simeq 2p.
\end{equation}
The corresponding fixed-target estimate grows only as the square root of the
beam momentum,
\begin{equation}
E_{\mathrm{cm}}^{\mathrm{fixed}}\sim \sqrt{mp}.
\end{equation}
The scaling, rather than the omitted numerical factor, is the point: doubling
the beam energy doubles the collider energy, but it increases fixed-target
center-of-mass energy only by a square root.

The invariant calculation makes the comparison precise. If \(P_1\) and
\(P_2\) are the incoming four-momenta, then
\begin{equation}
s=(P_1+P_2)^2,
\qquad E_{\mathrm{cm}}=\sqrt{s}.
\end{equation}
For equal counter-propagating beams of energy \(E\), the spatial momenta
cancel and \(\sqrt{s}=2E\). For a projectile of energy \(E\gg m\) striking an
equal-mass target at rest,
\begin{equation}
s=2m^2+2mE,
\qquad
\sqrt{s}\simeq \sqrt{2mE}
\qquad (c=1).
\end{equation}
This short invariant derivation supplies the factor suppressed in the spoken
estimate.\editorialnote{The derivation is included to make the lecture's
scaling argument reproducible.}

\begin{classroomqa}
\audiencequestion{Is there a fundamental reason a circular accelerator must
be so large? Could the same beam simply circulate through a smaller machine
more times?}
\lecturerresponse{Fields and materials impose limits. A particle receives a
finite increase in energy per metre, so a higher-energy machine needs more
effective accelerating length. A tight circle also demands greater transverse
acceleration, and an accelerated charge radiates. If the ring is too small,
the beam radiates energy as fast as the machine supplies it.}
\end{classroomqa}

The bending relation is approximately
\begin{equation}
p=qBR,
\end{equation}
so for fixed magnetic field \(B\), a larger momentum \(p\) requires a larger
radius \(R\). The radio-frequency cavities do the work that raises the beam
energy; the magnets bend and focus the beam. The broad lesson remains the one
given at the board: field strength and radiative losses prevent us from
compressing an arbitrarily energetic accelerator into a small ring.

Circular motion has acceleration even at constant speed because the velocity
changes direction. An accelerated charge emits synchrotron radiation. A
linear accelerator avoids the continuous bending loss, although the short
accelerating pushes themselves are not literally radiation-free. Circular
machines accept the bending cost in exchange for reusing the same apparatus on
every orbit.

\begin{classroomqa}
\audiencequestion{Why were electrons associated with the straight machine at
SLAC while circular high-energy colliders use protons or other hadrons?}
\lecturerresponse{Electrons are much lighter. At a given high energy their
Lorentz factor is much larger, and circular acceleration makes them radiate
far more strongly. Heavier protons can circle repeatedly, gaining energy on
each pass, before synchrotron losses become prohibitive.}
\end{classroomqa}

For fixed energy and ring radius, the radiated power scales very steeply with
the particle mass,
\begin{equation}
P_{\mathrm{sync}}\propto \frac{E^4}{m^4R^2}.
\end{equation}
This is why the electron--proton distinction is much sharper than the phrase
``the lighter particle moves faster'' might suggest: both beams can already be
extremely close to \(c\), but their Lorentz factors and radiation losses are
very different.\editorialnote{The bending relation and synchrotron scaling
make the lecture's qualitative accelerator argument precise; they were not
worked out on the board.}

\begin{classroomqa}
\audiencequestion{Why not accelerate hadrons in a straight line as well?}
\lecturerresponse{One can. The difficulty is length. In a ring the same
particle passes the accelerating structures many times; an equivalent linear
hadron machine would have to be enormously long. The useful compromise is a
large ring, large enough that proton radiation remains tolerable.}
\end{classroomqa}

These accelerator questions expose exactly why Newtonian rules are
insufficient. We now return to special relativity and build the language that
makes such comparisons systematic.

\section{The Garage and the Limousine}

Consider a garage about twelve feet long and a stretch limousine about twenty
feet long. The driver accelerates until the limousine is Lorentz-contracted,
hoping that it will fit. In the garage frame the moving limousine is shorter,
so it appears possible. In the limousine frame, however, the moving garage is
shorter, so it appears impossible. Can the whole limousine be inside the
garage?

Open both doors so that no collision with a wall distracts us. To say that the
limousine is inside means that its front and rear are inside
\emph{simultaneously}. That final word contains the entire puzzle.

Draw a spacetime diagram with \(t\) vertical and \(x\) horizontal. Light rays
form the \(45^\circ\) light cone in \(c=1\) units. The garage is at rest, so
the worldlines of its two doors are vertical. The front and rear of the
limousine trace two parallel sloping worldlines. The moving frame's \(t'\)-axis
is tilted toward the light cone, and its lines of simultaneity are parallel to
the tilted \(x'\)-axis.

\begin{figure}[H]
\centering
\includegraphics[width=0.82\textwidth]{lecture_05_figure_02.png}
\caption{The completed garage--limousine spacetime diagram. Vertical lines
bound the garage; sloping lines trace the two ends of the limousine. Horizontal
and tilted cuts represent the two frames' different simultaneity slices.}
\lectureframe{00:31:05}
\end{figure}

In the garage frame, a horizontal line can cross both limousine worldlines
between the two garage worldlines. At that garage time the rear has entered
and the front has not yet left. The contracted limousine is wholly inside.

In the limousine frame, the relevant simultaneous events lie on a tilted
line. Draw such a line through the event at which the rear enters. It meets
the front worldline after the front has already passed the far door. The
limousine is never wholly inside according to its own simultaneity. Both
statements are true because they compare different pairs of events.

\begin{classroomqa}
\audiencequestion{Do the two green worldlines and the horizontal segment
between them form a parallelogram? Is that black segment the length of the
vehicle?}
\lecturerresponse{The horizontal segment is the Lorentz-contracted length,
because its endpoints are simultaneous in the garage frame. The proper length
is obtained from endpoints simultaneous in the limousine frame, on a line
parallel to \(x'\). The two measurements use different cuts through the same
pair of worldlines.}
\end{classroomqa}

If \(L_0\) is the proper length, the horizontal garage-frame cut has length
\begin{equation}
L=\frac{L_0}{\gamma},
\qquad
\gamma=\frac{1}{\sqrt{1-v^2/c^2}}.
\end{equation}
The formula alone does not resolve the paradox. The diagram does, because it
shows which events were compared. Whenever a relativity puzzle seems to
produce two incompatible answers, look for an unspoken use of
``simultaneously.''

\section{Four Coordinates and Two Index Positions}

To move from puzzles to laws of motion, we need notation that remembers the
geometry. Package the coordinates as
\begin{equation}
(X^1,X^2,X^3,X^0)=(x,y,z,t),
\qquad X^\mu=(x,y,z,t),
\end{equation}
where \(\mu\) runs through \(1,2,3,0\). The ordering puts time last on the
page even though its index is \(0\), matching the convention used at the
board.

\begin{classroomqa}
\audiencequestion{Are we setting \(c=1\)? If ordinary units are restored,
should the time coordinate be \(ct\)?}
\lecturerresponse{Yes. We set \(c=1\) while developing the geometry and put
the factors of \(c\) back when dimensions or the nonrelativistic limit matter.
With conventional length units, the time coordinate is \(ct\).}
\end{classroomqa}

The same vector may be written with a lower index. With the sign convention
used here,
\begin{equation}
X_\mu=(-x,-y,-z,t).
\end{equation}
Thus lowering an index reverses the signs of the three spatial components and
leaves the time component unchanged. For example,
\begin{equation}
X^\mu=(3,2,4,10)
\quad\Longrightarrow\quad
X_\mu=(-3,-2,-4,10).
\end{equation}
These are not two spacetime points; they are two component forms of one
geometric object.

\begin{classroomqa}
\audiencequestion{Which components are covariant and contravariant, and are
the labels meant to be superscripts or subscripts?}
\lecturerresponse{Upper-index components are contravariant; lower-index
components are covariant. A useful memory aid from the room is ``co is low,
contra is high.'' Physicists often simply say upstairs and downstairs.}
\end{classroomqa}

The notation earns its keep in the invariant contraction
\begin{equation}
X_\mu X^\mu=t^2-x^2-y^2-z^2.
\label{eq:l5-interval}
\end{equation}
In differential form this is
\begin{equation}
d\tau^2=dt^2-dx^2-dy^2-dz^2
\qquad (c=1).
\label{eq:l5-proper-time}
\end{equation}
For timelike motion, \(d\tau\) is the time recorded by a clock moving along
the worldline.

Whenever the same index appears once upstairs and once downstairs, it is
understood to be summed. Einstein's convention turns
\begin{equation}
\sum_{\mu}X_\mu X^\mu
\end{equation}
into the shorter expression in Eq.~\eqref{eq:l5-interval}. This is not merely
typographical economy. It prevents the alternating signs of Minkowski
geometry from being rewritten by hand in every equation.

\begin{classroomqa}
\audiencequestion{Is \(\tau\) the standard symbol for proper time?}
\lecturerresponse{Yes. The symbol \(s\) is also sometimes used, but \(\tau\)
is a standard choice for proper time.}
\end{classroomqa}

\section{Rotations as the Easy Symmetry}

A Lorentz boost mixes space and time. Before studying that matrix, consider
the more familiar operation of rotating spatial axes. If the relative motion
points in an arbitrary direction, we can rotate the axes until that direction
becomes the new \(x\)-axis and then apply the known boost. Rotations are
therefore part of the construction, not an ornamental detour.

For a rotation through angle \(\theta\) about the \(z\)-axis,
\begin{align}
x'&=x\cos\theta+y\sin\theta,\nonumber\\
y'&=-x\sin\theta+y\cos\theta,\label{eq:l5-rotation}\\
z'&=z.\nonumber
\end{align}
At \(\theta=0\) these equations return \(x'=x\) and \(y'=y\); at
\(\theta=\pi/2\) they exchange the axes with the expected signs.

The transformation preserves Euclidean distance. Writing
\(c_\theta=\cos\theta\) and \(s_\theta=\sin\theta\),
\begin{align}
x'^2+y'^2
&=(xc_\theta+ys_\theta)^2+(-xs_\theta+yc_\theta)^2\nonumber\\
&=x^2(c_\theta^2+s_\theta^2)
  +y^2(s_\theta^2+c_\theta^2)\nonumber\\
&\quad+2xyc_\theta s_\theta-2xyc_\theta s_\theta\nonumber\\
&=x^2+y^2.
\end{align}
The cross terms cancel and \(c_\theta^2+s_\theta^2=1\). Since \(z'=z\),
the full length \(x^2+y^2+z^2\) is invariant.

The same equations become one matrix equation:
\begin{equation}
\begin{pmatrix}x'\\y'\\z'\end{pmatrix}
=
\underbrace{\begin{pmatrix}
\cos\theta&\sin\theta&0\\
-\sin\theta&\cos\theta&0\\
0&0&1
\end{pmatrix}}_{M}
\begin{pmatrix}x\\y\\z\end{pmatrix}.
\label{eq:l5-rotation-matrix}
\end{equation}
In index notation,
\begin{equation}
x'^i=M^i{}_{j}x^j,
\qquad i,j=1,2,3.
\end{equation}

\begin{figure}[H]
\centering
\includegraphics[width=0.82\textwidth]{lecture_05_figure_03.png}
\caption{The spatial rotation written both as a matrix equation and in index
notation.}
\lectureframe{01:00:40}
\end{figure}

\section{Lorentz Matrices and Composed Velocities}

For motion along \(x\), the Lorentz transformation in \(c=1\) units is
\begin{align}
x'&=\gamma_v(x-vt), & y'&=y,\nonumber\\
z'&=z, & t'&=\gamma_v(t-vx),
\end{align}
where
\begin{equation}
\gamma_v=\frac{1}{\sqrt{1-v^2}}.
\end{equation}
With the component order \((x,y,z,t)\), this becomes
\begin{equation}
\begin{pmatrix}x'\\y'\\z'\\t'\end{pmatrix}
=
\begin{pmatrix}
\gamma_v&0&0&-\gamma_vv\\
0&1&0&0\\
0&0&1&0\\
-\gamma_vv&0&0&\gamma_v
\end{pmatrix}
\begin{pmatrix}x\\y\\z\\t\end{pmatrix}.
\label{eq:l5-lorentz-matrix}
\end{equation}

\begin{figure}[H]
\centering
\includegraphics[width=0.82\textwidth]{lecture_05_figure_04.png}
\caption{The \(4\times4\) Lorentz matrix. Only the \(x\) and \(t\)
components mix; \(y\) and \(z\) remain unchanged.}
\lectureframe{01:04:40}
\end{figure}

Matrix multiplication now answers a physical question. Suppose one observer
moves with velocity \(v\) relative to us and a third moves with velocity \(u\)
relative to that observer, all along the same axis. Newtonian kinematics would
give \(w=u+v\), which would allow two velocities of \(3c/4\) to combine into
\(3c/2\). Relativity must give a different rule.

The active part of the boost is the \(2\times2\) matrix
\begin{equation}
B(v)=\gamma_v
\begin{pmatrix}1&-v\\-v&1\end{pmatrix}.
\end{equation}
Composing the two boosts gives
\begin{align}
B(u)B(v)
&=\gamma_u\gamma_v
\begin{pmatrix}1&-u\\-u&1\end{pmatrix}
\begin{pmatrix}1&-v\\-v&1\end{pmatrix}\nonumber\\
&=\gamma_u\gamma_v
\begin{pmatrix}
1+uv&-(u+v)\\
-(u+v)&1+uv
\end{pmatrix}.
\label{eq:l5-composed-boost}
\end{align}
The ratio of an off-diagonal entry to a diagonal entry identifies the velocity
of the resulting boost:
\begin{equation}
w=\frac{u+v}{1+uv}
\qquad(c=1).
\end{equation}
Restoring dimensions,
\begin{equation}
w=\frac{u+v}{1+uv/c^2}.
\label{eq:l5-velocity-addition}
\end{equation}

\begin{figure}[H]
\centering
\includegraphics[width=0.82\textwidth]{lecture_05_figure_05.png}
\caption{The composed-boost matrix and the velocity
\(w=(u+v)/(1+uv)\) in natural units.}
\lectureframe{01:15:35}
\end{figure}

If \(u\) and \(v\) are each close to \(c\), the numerator approaches \(2c\)
while the denominator approaches \(2\), so \(w\) approaches \(c\), not
\(2c\). More generally, Eq.~\eqref{eq:l5-velocity-addition} maps subluminal
inputs to a subluminal result. Matrix multiplication has encoded the speed
limit automatically.

\section{Four-Vectors and the Failure of Ordinary Velocity}

A four-vector is any object whose four components transform under rotations
and Lorentz boosts in the same way as the spacetime coordinates. If \(V^\mu\)
is a four-vector, then with our sign convention
\begin{equation}
V_\mu=(-V^1,-V^2,-V^3,V^0),
\end{equation}
and its invariant norm is
\begin{equation}
V_\mu V^\mu=(V^0)^2-(V^1)^2-(V^2)^2-(V^3)^2.
\end{equation}
For two four-vectors \(V^\mu\) and \(W^\mu\), the contraction
\begin{equation}
W_\mu V^\mu
=W^0V^0-W^1V^1-W^2V^2-W^3V^3
\end{equation}
is likewise invariant. The components change between frames; the contraction
does not.

Ordinary velocity has only three components,
\begin{equation}
\mathbf v=\left(\frac{dx}{dt},\frac{dy}{dt},\frac{dz}{dt}\right).
\end{equation}
Appending \(dt/dt=1\) does not make it a useful four-vector. The fourth entry
would be trivial, and the resulting object would not transform linearly under
Lorentz transformations. Coordinate time \(t\) changes from one frame to
another, so differentiating by \(t\) spoils the desired transformation law.

Proper time provides the cure. Along a worldline define
\begin{equation}
u^\mu=\frac{dX^\mu}{d\tau}.
\label{eq:l5-four-velocity}
\end{equation}
The differential \(dX^\mu\) is a four-vector and \(d\tau\) is invariant, so
their ratio is a four-vector. The adjective ``proper'' refers to the clock in
the denominator.

The chain rule relates this object to ordinary velocity:
\begin{equation}
\frac{dx}{d\tau}=\frac{dx}{dt}\frac{dt}{d\tau},
\end{equation}
and similarly for \(y\) and \(z\). From
Eq.~\eqref{eq:l5-proper-time},
\begin{align}
\left(\frac{d\tau}{dt}\right)^2
&=1-\left(\frac{dx}{dt}\right)^2
    -\left(\frac{dy}{dt}\right)^2
    -\left(\frac{dz}{dt}\right)^2\nonumber\\
&=1-v^2
\qquad(c=1),
\end{align}
and therefore
\begin{equation}
\frac{dt}{d\tau}=\gamma_v=\frac{1}{\sqrt{1-v^2}}.
\end{equation}
The four components are
\begin{equation}
u^i=\gamma_v v^i,
\qquad
u^0=\gamma_v
\qquad(c=1).
\label{eq:l5-four-velocity-components}
\end{equation}
At low speed, \(u^i\simeq v^i\) and \(u^0\simeq1\). Near the speed of light,
\(\gamma_v\) grows without bound and proper velocity differs sharply from
ordinary velocity.

\begin{figure}[H]
\centering
\includegraphics[width=0.82\textwidth]{lecture_05_figure_06.png}
\caption{The four components of proper velocity expressed through ordinary
velocity and the Lorentz factor.}
\lectureframe{01:34:40}
\end{figure}

Restoring conventional units and using \(X^0=ct\) gives the standard form
\begin{equation}
U^\mu=(\gamma_v\mathbf v,\gamma_v c)
\end{equation}
in the spatial-first ordering used here. Its invariant norm is
\begin{equation}
U_\mu U^\mu=c^2.
\end{equation}
Thus every massive particle has a four-velocity of the same invariant
magnitude even though its spatial velocity depends on the frame.

\section{The Tangent to a Worldline}

There is a familiar Euclidean model for Eq.~\eqref{eq:l5-four-velocity}.
Along a curve in ordinary three-dimensional space,
\begin{equation}
ds^2=dx^2+dy^2+dz^2.
\end{equation}
The differential displacement \((dx,dy,dz)\) points along the curve but is
not a unit vector. Dividing by arc length gives
\begin{equation}
\mathbf T=
\left(\frac{dx}{ds},\frac{dy}{ds},\frac{dz}{ds}\right),
\qquad
\mathbf T\cdot\mathbf T=1.
\end{equation}
The relativistic construction is the same idea with proper time in place of
Euclidean arc length. The vector \(u^\mu=dX^\mu/d\tau\) is tangent to the
worldline.

\begin{figure}[H]
\centering
\includegraphics[width=0.82\textwidth]{lecture_05_figure_07.png}
\caption{The Euclidean tangent-vector construction: divide the differential
displacement by the arc-length element \(ds\). Proper time plays the analogous
role along a timelike worldline.}
\lectureframe{01:39:30}
\end{figure}

\begin{classroomqa}
\audiencequestion{If \(c=1\), are proper time and proper distance the same?}
\lecturerresponse{Setting \(c=1\) puts time and distance in the same units; it
does not erase the minus signs of spacetime geometry. In the Euclidean analogy,
\(s\) is ordinary distance measured along a spatial curve. Along a timelike
worldline, \(\tau\) is the elapsed time measured by a clock. The constructions
are analogous, but the geometries are different.}
\end{classroomqa}

Differentiating once more defines proper acceleration,
\begin{equation}
a^\mu=\frac{du^\mu}{d\tau}.
\end{equation}
This is the quantity that must replace \(d^2\mathbf x/dt^2\) when Newton's
second law is made relativistic. Since \(u_\mu u^\mu\) is constant,
\begin{equation}
\frac{d}{d\tau}(u_\mu u^\mu)=2u_\mu a^\mu=0,
\end{equation}
so four-acceleration is Minkowski-orthogonal to four-velocity.\editorialnote{The
orthogonality follows immediately from the invariant norm and is included as
the next mathematical consequence of the definition.}

\section{Four-Momentum and Energy}

Momentum is the final example because it is the conserved quantity needed for
collisions and decays. In modern terminology the mass \(m\) is an intrinsic
property of the particle---what older books called the rest mass---rather than
a quantity that changes with velocity. Relativistic momentum is obtained by
multiplying that fixed mass by four-velocity:
\begin{equation}
p^\mu=mu^\mu.
\label{eq:l5-four-momentum}
\end{equation}
In \(c=1\) units,
\begin{equation}
p^i=\gamma_v m v^i,
\qquad
p^0=\gamma_v m.
\end{equation}
At small velocity \(\gamma_v\simeq1\), so the spatial components reduce to
the Newtonian momentum \(m\mathbf v\). At high velocity they grow more rapidly
than \(m\mathbf v\).

The time component is energy. With conventional units and \(X^0=ct\),
\begin{equation}
P^\mu=(\mathbf p,E/c)
=(\gamma_v m\mathbf v,\gamma_v mc),
\end{equation}
and hence
\begin{equation}
E=\gamma_v mc^2
=mc^2\left(1-\frac{v^2}{c^2}\right)^{-1/2}.
\label{eq:l5-relativistic-energy}
\end{equation}

To recover ordinary mechanics, let
\(\epsilon=v^2/c^2\ll1\). The binomial expansion gives
\begin{equation}
(1-\epsilon)^p=1-p\epsilon+O(\epsilon^2).
\end{equation}
Taking \(p=-1/2\),
\begin{align}
E
&=mc^2\left(1-\frac{v^2}{c^2}\right)^{-1/2}\nonumber\\
&=mc^2\left[1+\frac12\frac{v^2}{c^2}
  +O\!\left(\frac{v^4}{c^4}\right)\right]\nonumber\\
&=mc^2+\frac12mv^2+O\!\left(\frac{mv^4}{c^2}\right).
\label{eq:l5-energy-expansion}
\end{align}

\begin{figure}[H]
\centering
\includegraphics[width=0.82\textwidth]{lecture_05_figure_08.png}
\caption{The low-velocity expansion of relativistic energy, revealing rest
energy and Newtonian kinetic energy.}
\lectureframe{01:51:30}
\end{figure}

The first term, \(mc^2\), is present even at rest. In a Newtonian problem with
a fixed collection of particles and unchanging masses, that term adds the
same constant before and after the motion, so only \(\tfrac12mv^2\) affects
the familiar bookkeeping. Relativity reveals that the discarded constant was
physical energy all along.

Energy and momentum are therefore not unrelated conservation laws. They are
the four components of one conserved vector. A Lorentz boost mixes them: what
one observer calls energy contributes partly to another observer's momentum,
and conversely. Nevertheless, for an isolated system the total four-momentum
is the same before and after every collision or decay in every inertial frame.

This is the threshold of relativistic mechanics. The next tasks are to express
force through proper acceleration and to understand relativistic wave
equations, especially the equations governing electromagnetism and light.
